\documentclass[aspectratio=169,12pt]{beamer}

% 中文支援
\usepackage{xeCJK}
\usepackage{fontspec}
\setCJKmainfont{Noto Serif CJK TC}
\setCJKsansfont{Noto Sans CJK TC}

% 其他套件
\usepackage{graphicx}
\usepackage{booktabs}
\usepackage{array}
\usepackage{longtable}
\usepackage{tikz}
\usetikzlibrary{shapes.geometric, arrows}

% Beamer主題設定
\usetheme{Madrid}
\usecolortheme{default}

% 自訂顏色
\definecolor{emergency_red}{RGB}{186,24,27}
\definecolor{emergency_blue}{RGB}{0,114,188}
\definecolor{emergency_gray}{RGB}{120,120,120}
\definecolor{light_blue}{RGB}{173,216,230}
\definecolor{light_green}{RGB}{144,238,144}

\setbeamercolor{structure}{fg=emergency_blue}
\setbeamercolor{title}{fg=white}
\setbeamercolor{frametitle}{fg=white}
\setbeamercolor{block title}{bg=emergency_blue,fg=white}
\setbeamercolor{block body}{bg=light_blue!30}

% 標題資訊
\title[CCC委員會章程]{內科部臨床能力委員會工作計畫章程}
\subtitle{Clinical Competency Committee Charter}
\author{○○醫院內科部}
\date{版本 1.4 \\ \today}
\institute{內科部教學委員會}

\begin{document}

% 標題頁
\begin{frame}
\titlepage
\end{frame}

% 目錄
\begin{frame}{簡報大綱}
\tableofcontents
\end{frame}

\section{章程總則}

\begin{frame}{設立依據}
\begin{block}{法規基礎}
本臨床能力委員會（CCC）依據以下法規設立：
\end{block}
\begin{itemize}
    \item 衛生福利部「專科醫師訓練課程基準」
    \item 「畢業後一般醫學訓練計畫基準」
    \item 醫院評鑑暨醫療品質策進會相關標準
    \item 本院醫學教育委員會授權
    \item 內科部教學訓練規章
    \item 「住院醫師勞動權益保障及工作時間指引」
\end{itemize}
\end{frame}

\begin{frame}{委員會宗旨}
\begin{alertblock}{核心理念}
因應住院醫師納入勞基法後工時縮減的挑戰，建立以勝任能力為本（CBME）的系統性教學評量機制
\end{alertblock}

\begin{block}{教育模式轉型}
\begin{itemize}
    \item 從「以時間為本」轉向「以能力為本」
    \item 確保各層級醫事人員達到專業標準
    \item 維護醫療品質與病患安全
    \item 善盡醫學教育之社會責任
\end{itemize}
\end{block}
\end{frame}

\begin{frame}{適用範圍}
\begin{columns}[T]
\column{0.5\textwidth}
\begin{block}{主要對象}
\begin{itemize}
    \item 醫學系學生
    \item PGY醫師
    \item 住院醫師（R1-R3）
\end{itemize}
\end{block}

\column{0.5\textwidth}
\begin{block}{其他對象}
\begin{itemize}
    \item 專科護理師
    \item 總住院醫師
    \item 其他經委員會決議納入評核之醫事人員
\end{itemize}
\end{block}
\end{columns}
\end{frame}

\begin{frame}{核心理念}
\begin{block}{ACGME NAS架構}
本委員會採用美國ACGME「下一代評鑑系統」結合台灣本土特色
\end{block}

\begin{itemize}
    \item \textbf{結構嚴謹}：標準化評核流程與決議程序
    \item \textbf{證據導向}：多元評量工具整合運用
    \item \textbf{階段發展}：依訓練階段設定能力里程碑
    \item \textbf{核心勝任能力}：以ACGME六大核心能力為架構
    \item \textbf{完善回饋}：即時性、明確性、個人化回饋
\end{itemize}
\end{frame}

\begin{frame}{委員會功能定位}
\begin{columns}[T]
\column{0.5\textwidth}
\begin{block}{能力發展模型 ✓}
\begin{itemize}
    \item 系統性詮釋評量工具
    \item 里程碑評量地圖
    \item 整合多項評量資料
    \item 形成性回饋
    \item 個人化學習指引
\end{itemize}
\end{block}

\column{0.5\textwidth}
\begin{alertblock}{避免問題辨認模型 ✗}
\begin{itemize}
    \item 避免「舉紅旗」模式
    \item 不僅識別困境者
    \item 為每位受訓者提供建設性指導
\end{itemize}
\end{alertblock}
\end{columns}
\end{frame}

\section{委員會組織架構}

\begin{frame}{委員會組成}
\begin{block}{委員總數}
15-20名委員，確保代表性與專業性
\end{block}

\begin{table}
\footnotesize
\begin{tabular}{lcc}
\toprule
\textbf{委員類別} & \textbf{人數} & \textbf{任期} \\
\midrule
委員會主席 & 1名 & 2年 \\
次專科主任代表 & 6-8名 & 1年 \\
醫事人員代表 & 4-6名 & 1年 \\
外部委員 & 1-2名 & 2年 \\
\bottomrule
\end{tabular}
\end{table}
\end{frame}

\begin{frame}{委員資格}
\begin{columns}[T]
\column{0.5\textwidth}
\begin{block}{基本資格}
\begin{itemize}
    \item 相關專科證書
    \item 本院服務滿2年
    \item 教學或管理經驗
    \item 無重大糾紛記錄
    \item 完成CCC培訓
\end{itemize}
\end{block}

\column{0.5\textwidth}
\begin{block}{專業資格}
\begin{itemize}
    \item 主治醫師：教學年資3年以上
    \item 醫事人員：年資5年以上
    \item 外部委員：資深或教育專家
\end{itemize}
\end{block}
\end{columns}
\end{frame}

\begin{frame}{利益迴避原則}
\begin{alertblock}{應迴避情形}
\begin{itemize}
    \item 受評核者為其直接指導學生
    \item 與受評核者有直系血親、姻親關係
    \item 與受評核者有重大利益關係
    \item 其他可能影響公正性之情形
\end{itemize}
\end{alertblock}

\begin{exampleblock}{重點提醒}
指導老師在評核其直接指導學生時，可提供意見但不得參與表決
\end{exampleblock}
\end{frame}

\section{委員會職責與權限}

\begin{frame}{主要職責（1/2）}
\begin{block}{能力評核認定}
\begin{itemize}
    \item 制定並修訂臨床能力評核標準
    \item 執行定期能力評核作業
    \item 審核學員能力進階申請
    \item 認定學員是否達到畢業標準
\end{itemize}
\end{block}

\begin{block}{品質監控}
\begin{itemize}
    \item 監控教學訓練品質
    \item 分析學員表現趨勢
    \item 識別訓練課程缺失
    \item 提出改善建議方案
\end{itemize}
\end{block}
\end{frame}

\begin{frame}{主要職責（2/2）}
\begin{block}{政策制定}
\begin{itemize}
    \item 制定內科部教學政策
    \item 修訂訓練課程基準
    \item 建立評核作業流程
    \item 訂定獎懲標準
\end{itemize}
\end{block}
\end{frame}

\begin{frame}{評核對象與標準：醫學生}
\begin{block}{評核週期}
每1-3週評核一次
\end{block}

\begin{itemize}
    \item \textbf{第1次評核（1週）}：基本知識與臨床技能
    \item \textbf{第2次評核（2-3週）}：臨床判斷能力
\end{itemize}

\begin{exampleblock}{重點能力}
\begin{itemize}
    \item 病史詢問與理學檢查
    \item 臨床推理能力
    \item 醫病溝通技巧
\end{itemize}
\end{exampleblock}
\end{frame}

\begin{frame}{評核對象與標準：PGY醫師}
\begin{block}{評核週期}
每3個月評核一次，共4次評核
\end{block}

\begin{table}
\footnotesize
\begin{tabular}{ll}
\toprule
\textbf{評核時程} & \textbf{評核重點} \\
\midrule
第1次（3個月） & 基礎臨床能力 \\
第2次（6個月） & 獨立處置能力 \\
第3次（9個月） & 複雜病例處理 \\
第4次（12個月） & 整合能力評估 \\
\bottomrule
\end{tabular}
\end{table}
\end{frame}

\begin{frame}{評核對象與標準：住院醫師}
\begin{block}{評核週期}
每3-6個月評核一次
\end{block}

\begin{description}
    \item[R1] 基礎臨床照護能力、團隊合作
    \item[R2] 獨立診療能力、進階處置技能
    \item[R3] 次專科訓練、教學能力、領導能力
\end{description}

\vspace{0.5em}
\begin{alertblock}{畢業標準}
須達到所有核心EPA的Level 4（獨立執行）水準
\end{alertblock}
\end{frame}

\section{評核流程與機制}

\begin{frame}{會議召開規範}
\begin{block}{定期會議}
\begin{itemize}
    \item 每季召開一次（每3個月）
    \item 至少提前7日通知
    \item 應有2/3以上委員出席
\end{itemize}
\end{block}

\begin{block}{臨時會議}
\begin{itemize}
    \item 主席或1/3以上委員提議
    \item 處理緊急或重大案件
    \item 應有1/2以上委員出席
\end{itemize}
\end{block}
\end{frame}

\begin{frame}{評核程序（1/3）}
\begin{enumerate}
    \item \textbf{資料準備階段}
    \begin{itemize}
        \item 秘書彙整評量資料
        \item 製作評核摘要報告
        \item 提前3日提供委員
    \end{itemize}
    
    \item \textbf{會議評核階段}
    \begin{itemize}
        \item 主席主持會議
        \item 報告評核案件
        \item 委員審查討論
    \end{itemize}
\end{enumerate}
\end{frame}

\begin{frame}{評核程序（2/3）}
\begin{enumerate}
\setcounter{enumi}{2}
    \item \textbf{決議表決階段}
    \begin{itemize}
        \item 委員投票表決
        \item 需2/3以上同意通過
        \item 記錄決議結果
    \end{itemize}
    
    \item \textbf{回饋通知階段}
    \begin{itemize}
        \item 3個工作天內通知
        \item 提供書面回饋
        \item 安排面談說明
    \end{itemize}
\end{enumerate}
\end{frame}

\begin{frame}{評核程序（3/3）}
\begin{enumerate}
\setcounter{enumi}{4}
    \item \textbf{追蹤執行階段}
    \begin{itemize}
        \item 監督改善計畫執行
        \item 定期檢視進度
        \item 必要時調整計畫
    \end{itemize}
\end{enumerate}

\begin{alertblock}{重要原則}
整個流程應遵循保密原則、公平原則、及時原則
\end{alertblock}
\end{frame}

\begin{frame}{評量工具整合}
\begin{block}{多元評量來源}
\begin{itemize}
    \item Mini-CEX（臨床技能評估）
    \item DOPS（操作型評量）
    \item CbD（病例討論）
    \item 360度評估（多面向回饋）
    \item 工作量統計
    \item 臨床表現記錄
\end{itemize}
\end{block}
\end{frame}

\begin{frame}{決議機制}
\begin{block}{表決方式}
\begin{itemize}
    \item 一般決議：出席委員1/2以上同意
    \item 重大決議：出席委員2/3以上同意
    \item 主席有最終決定權
\end{itemize}
\end{block}

\begin{exampleblock}{重大決議包括}
\begin{itemize}
    \item 延長訓練期
    \item 不予晉級
    \item 建議退訓
\end{itemize}
\end{exampleblock}
\end{frame}

\begin{frame}{評核結果分類}
\begin{table}
\footnotesize
\begin{tabular}{lp{7cm}}
\toprule
\textbf{結果} & \textbf{說明} \\
\midrule
通過 & 達到該階段能力標準，可繼續訓練 \\
有條件通過 & 大致達標但需加強特定項目 \\
需改善 & 未達標準，須執行改善計畫並重新評核 \\
不通過 & 嚴重未達標準，需延長訓練或其他處置 \\
\bottomrule
\end{tabular}
\end{table}
\end{frame}

\section{申訴與救濟程序}

\begin{frame}{申訴權利}
\begin{alertblock}{申訴期限}
受評核者對評核結果不服時，得於收到通知後\textbf{15日內}提出申訴
\end{alertblock}
\end{frame}

\begin{frame}{申訴程序}
\begin{block}{第一階段：內部申訴}
\begin{enumerate}
    \item 向委員會秘書提出書面申訴
    \item 委員會於30日內重新審查
    \item 申訴人得要求列席說明
    \item 委員會作成決議並書面回覆
\end{enumerate}
\end{block}

\begin{block}{第二階段：院級申訴}
\begin{enumerate}
    \item 對第一階段結果不服，向院級申訴委員會申訴
    \item 院級申訴委員會於45日內作成最終決議
    \item 申訴結果為最終決議，不得再行申訴
\end{enumerate}
\end{block}
\end{frame}

\begin{frame}{申訴期間權益保障}
\begin{exampleblock}{受保障權益}
\begin{itemize}
    \item 申訴期間得繼續原訓練計畫
    \item 不得因申訴而受到不當對待
    \item 申訴資料保密處理
    \item 提供適當的法律諮詢協助
\end{itemize}
\end{exampleblock}
\end{frame}

\section{爭議處理}

\begin{frame}{爭議處理程序}
\begin{block}{內部爭議處理}
\begin{enumerate}
    \item 委員間意見分歧：主席協調，必要時表決
    \item 評估結果爭議：重新檢視證據，再次討論
    \item 程序問題：參照章程規定，法制組諮詢
    \item 倫理疑慮：醫學倫理委員會介入處理
\end{enumerate}
\end{block}
\end{frame}

\begin{frame}{外部申訴處理}
\begin{block}{處理原則}
\begin{enumerate}
    \item 受訓者申訴：啟動正式申訴程序
    \item 指導老師質疑：提供詳細說明和證據
    \item 外部關切：透明化處理過程和結果
    \item 媒體關注：統一由發言人對外說明
\end{enumerate}
\end{block}
\end{frame}

\section{總結}

\begin{frame}{重要原則回顧}
\begin{block}{核心價值}
\begin{itemize}
    \item 以能力發展為導向，非問題識別
    \item 提供每位受訓者個人化回饋
    \item 建立系統性、證據導向的評核機制
    \item 確保評核的公平性與客觀性
    \item 維護受訓者權益與申訴管道
\end{itemize}
\end{block}
\end{frame}

\begin{frame}{章程效力}
\begin{alertblock}{法律效力}
本章程經內科部主任指導，由內科部修訂、經教學部同意備查，所有相關人員均須遵守。違反者將依相關規定處理。
\end{alertblock}

\vspace{1em}

\begin{center}
\textbf{版本 1.4}\\
生效日期：\today
\end{center}
\end{frame}

\begin{frame}[plain]
\begin{center}
\Huge 謝謝聆聽\\[1em]
\Large 敬請指教
\end{center}
\end{frame}

\end{document}